\section{Introduction}
Formula Student (FS) is an international engineering competition in which university teams design, build and race a single-seat car; in the driverless (DV) class, the same car must complete the dynamic events autonomously. Every dynamic event is driven on a track that exists only as a pattern of color coded cones on an open area of tarmac. Under the 2026 driverless specification, the left border is marked with small blue cones, the right border with small yellow cones, entry and exit lanes with small orange cones, and big orange cones are placed before and after the start, finish and timekeeping lines, with at most 5 m between consecutive cones in the driving direction and shorter spacing in corners \cite{fsds2026}. Until to now, these layouts are reproduced many times per season by the team during test days: acceleration and skidpad layouts for vehicle tuning, and autocross and trackdrive circuits of a few hundred metres, which already need well over a hundred cones per lap.

These tracks are set out by hand: a plan drawn on paper or in a generated track is transferred to the ground with measuring tapes and trigonometry, which needs several people, takes time away from testing and provides no check of the result: a misplaced or wrong-coloured cone is only corrected if someone notices it. For driverless cars, the consequences go beyond inconvenience. Their perception and mapping pipelines use the cones as the landmarks that define the track, and path planning and control are built on the map estimated from them \cite{kabzan2020}. A layout that differs from the intended one changes the problem the car is solving and makes results from different test days hard to compare. Results are also only comparable if each session is linked to the car that ran, which is today recorded by hand, if at all. Surveying instruments such as total stations and real-time kinematic (RTK) global navigation satellite system (GNSS) receivers reach centimetre accuracy \cite{jiang2023}, but they are dedicated equipment that student teams rarely own or are trained to use. None of these options links the plan, the placement on the field and a record of what was actually built.

Augmented reality (AR) combines virtual content with the user's view of the real world, registered in three dimensions and interactive in real time \cite{azuma1997}. Mobile AR frameworks on consumer smartphones, such as Google's ARCore, estimate the pose of the device by combining visual feature points from the camera with inertial measurements \cite{arcore}. This makes it possible to show each planned cone, with its colour, at its position on the ground through a smartphone, so that a single team member can set out the track alone. Compared with the conventional alternatives (a tape measure, or a 2D map showing the phone's GNSS position), AR guidance removes the measuring step, shows the instruction where the cone must be placed, and allows the same camera to verify the placement. Evidence from construction point layout shows both the potential and the main limitation: with a head-mounted AR device registered to a paper marker, an accuracy of about 5 cm was reported, which became worse further from the marker \cite{chalhoub2019}. Their indoor test area was much smaller than an FS test site, so this accuracy cannot be assumed to hold over a full track.

The proposed system, TrackSeed, combines the contributions of the two course units around a shared model of the track. Virtual and Augmented Reality (VAR) provides the AR guidance aligned to the real ground and a digital twin (DT) of the track: a live virtual model that is updated automatically with every cone confirmed by the phone, and whose edits on the dashboard are sent back to the phone as updated guidance. A full DT exchanges data automatically in both directions; because a person still has to move the physical cone, TrackSeed sits between a digital shadow and a full twin \cite{kritzinger2018}. Emerging and Intelligent Technologies (EIT) provides the cloud platform that stores track plans and placement sessions, real-time synchronisation between the phone and a web dashboard, remote configuration of the app, and on-device intelligence that checks each placed cone with the camera. In this first version, one person places the cones with one phone; collaborative placement with several phones is left for future work.

This paper reviews the state of the art that TrackSeed builds on, with three objectives: (i) to establish what accuracy mobile AR achieves for on-site layout and guidance, and which alignment approaches limit drift over large outdoor areas; (ii) to compare architectures for keeping a DT synchronised in real time between mobile and web clients; and (iii) to derive and justify the technology choices and validation plan of the proposal. Section II reviews the literature by theme, Section III analyses the candidate technologies, Section IV presents the proposed solution and its validation plan, Section V discusses feasibility and limitations, and Section VI concludes.

\subsection{Scope and literature search}
The literature was searched in October 2026 in Google Scholar and in publishers' digital libraries, assisted by an AI tool (see the AI use declaration) and by following the references and citations of the most relevant studies. Search strings were built from the main concepts of the problem (AR, on-site layout, mobile tracking accuracy and digital twins) and their synonyms, for example (``augmented reality'' OR ``mixed reality'') AND (``layout'' OR ``setout'' OR ``stakeout'') AND (``accuracy'' OR ``error''), and (``ARCore'' OR ``mobile AR'') AND (``drift'' OR ``accuracy'') AND ``outdoor''. A source was kept only if it was directly relevant to the problem, reported evidence that could be compared, came from a credible venue, was recent enough to reflect current devices, and helped compare alternatives or justify a design decision; peer-reviewed papers published from 2015 onwards were preferred. Foundational works were kept for definitions, and official documentation and the FS rules were used only for tool capabilities and domain requirements, not as evidence of performance. Since almost no studies address track setout in motorsport, the review relies mainly on adjacent domains, construction layout in particular, whose tasks, distances and accuracy requirements differ from those of an FS test site.
